\documentclass[12pt,a4paper]{article}

% ------------------------------------------------
% PAQUETES
% ------------------------------------------------

\usepackage[utf8]{inputenc}
\usepackage[T1]{fontenc}
\usepackage[spanish]{babel}

\usepackage{geometry}
\geometry{
  left=2.5cm,
  right=2.5cm,
  top=2.5cm,
  bottom=2.5cm
}

\usepackage{array}
\usepackage{tabularx}
\usepackage{booktabs}
\usepackage{longtable}
\usepackage{enumitem}
\usepackage{xcolor}
\usepackage{hyperref}
\usepackage{graphicx}
\usepackage{listings}
\usepackage{fancyhdr}
\usepackage{titlesec}
\usepackage{tcolorbox}
\usepackage{multirow}

\tcbuselibrary{breakable, skins}

\hypersetup{
  colorlinks=true,
  linkcolor=blue,
  urlcolor=blue
}

% Estilo de código
\lstset{
  basicstyle=\ttfamily\small,
  breaklines=true,
  frame=single,
  backgroundcolor=\color{gray!8},
  keywordstyle=\color{blue}\bfseries,
  commentstyle=\color{green!50!black},
  stringstyle=\color{red!70!black},
  showstringspaces=false,
  numbers=left,
  numberstyle=\tiny\color{gray},
  numbersep=5pt
}

% Cajas informativas
\newtcolorbox{infobox}[1][]{
  colback=blue!5,
  colframe=blue!40,
  fonttitle=\bfseries,
  title=#1,
  breakable
}

\newtcolorbox{warningbox}[1][]{
  colback=orange!8,
  colframe=orange!60,
  fonttitle=\bfseries,
  title=#1,
  breakable
}

% Encabezado y pie de página
\pagestyle{fancy}
\fancyhf{}
\rhead{Ingeniería de Software}
\lhead{Hito 3 --- SkyConnect Airlines}
\rfoot{Página \thepage}
\renewcommand{\headrulewidth}{0.4pt}

% ------------------------------------------------
% DATOS DEL DOCUMENTO
% ------------------------------------------------

\title{
  \textbf{HITO 3}\\[0.4cm]
  \large Optimización y Escalabilidad de una Plataforma\\
  de Venta de Pasajes Aéreos bajo Alta Concurrencia
}

\author{Curso: Ingeniería de Software}
\date{}

% ================================================
\begin{document}
% ================================================

\maketitle
\thispagestyle{fancy}

% ------------------------------------------------
\section{Antecedentes}
% ------------------------------------------------

Durante eventos masivos de comercio electrónico como \textit{Cyber Day} o \textit{Black Friday}, las plataformas de venta de pasajes aéreos experimentan incrementos significativos en la cantidad de usuarios concurrentes. Esta situación genera desafíos importantes en rendimiento, disponibilidad y escalabilidad de los sistemas informáticos.

La empresa ficticia \textbf{SkyConnect Airlines} ha detectado que su plataforma presenta tiempos de respuesta inaceptables cuando miles de usuarios intentan consultar vuelos simultáneamente, afectando la experiencia de compra y provocando pérdidas económicas directas.

Cada equipo recibirá un \textbf{sistema base completamente funcional y dockerizado}. El trabajo consiste en analizar sus limitaciones bajo carga, identificar cuellos de botella y aplicar optimizaciones medibles, documentando cada decisión técnica y económica.

% ------------------------------------------------
\section{Objetivo General}
% ------------------------------------------------

Analizar, medir y optimizar el rendimiento de una plataforma de venta de pasajes aéreos bajo condiciones de alta concurrencia, justificando técnica y económicamente las mejoras implementadas.

% ------------------------------------------------
\section{Objetivos Específicos}
% ------------------------------------------------

Al finalizar la actividad los estudiantes serán capaces de:

\begin{itemize}
  \item Identificar y cuantificar cuellos de botella presentes en una arquitectura web real.
  \item Implementar mecanismos de caché utilizando Redis con estrategias de invalidación adecuadas.
  \item Aplicar herramientas de observabilidad (Prometheus y Grafana) para medir el impacto de las mejoras.
  \item Diseñar y ejecutar pruebas de carga reproducibles con Locust.
  \item Comparar métricas objetivas antes y después de cada optimización.
  \item Justificar técnica y económicamente las decisiones de arquitectura adoptadas.
\end{itemize}

% ------------------------------------------------
\section{Sistema Base Entregado}
% ------------------------------------------------

Cada equipo recibirá una aplicación completamente funcional y dockerizada. \textbf{No se requiere implementar funcionalidad nueva de negocio}: el trabajo se centra en medir, optimizar y documentar.

\subsection{Stack tecnológico}

\begin{itemize}
  \item \textbf{Frontend:} React + Axios + Bootstrap
  \item \textbf{Backend:} FastAPI (Python) + SQLAlchemy (ORM)
  \item \textbf{Base de datos:} PostgreSQL
  \item \textbf{Orquestación:} Docker Compose
\end{itemize}

\subsection{Estructura del proyecto}

\begin{lstlisting}[language=bash, caption={Estructura de directorios del sistema base}]
skyconnect/
  backend/
    app/
      main.py          # Punto de entrada FastAPI
      models.py        # Modelos SQLAlchemy
      schemas.py       # Esquemas Pydantic
      database.py      # Conexion a PostgreSQL
      routers/
        flights.py     # Endpoint GET /flights
        routes.py      # Endpoint GET /routes
        aircraft.py    # Endpoint GET /aircraft
    requirements.txt
    Dockerfile
  frontend/
    src/
      App.jsx
      components/
        FlightSearch.jsx
        FlightList.jsx
    Dockerfile
  db/
    init.sql           # Script de inicializacion y seed de datos
  docker-compose.yml
  locust/
    locustfile.py      # Script base de pruebas de carga
\end{lstlisting}

\subsection{Levantar el sistema base}

\begin{lstlisting}[language=bash, caption={Comandos para levantar el entorno}]
# Clonar el repositorio del sistema base
git clone <url-repositorio-base>
cd skyconnect

# Levantar todos los servicios
docker compose up --build

# Verificar que los servicios estan activos
docker compose ps

# Acceder a la aplicacion
# Frontend:  http://localhost:3000
# Backend:   http://localhost:8000
# API Docs:  http://localhost:8000/docs
\end{lstlisting}

\subsection{Arquitectura inicial}

\begin{verbatim}
+---------------+         HTTP          +---------------+
|               |  ------------------>  |               |
|    React      |                       |    FastAPI    |
|  (Frontend)   |  <------------------  |   (Backend)   |
|               |       JSON            |               |
+---------------+                       +-------+-------+
                                                |
                                         SQL queries
                                                |
                                        +-------v-------+
                                        |               |
                                        |  PostgreSQL   |
                                        |               |
                                        +---------------+
\end{verbatim}

\begin{warningbox}[Problema observable]
Cada request al endpoint \texttt{GET /flights} ejecuta múltiples \texttt{JOIN} sobre PostgreSQL, incluyendo cálculo dinámico de asientos disponibles. Con alta concurrencia, la base de datos se convierte en el cuello de botella principal, degradando el tiempo de respuesta significativamente.
\end{warningbox}

\subsection{Arquitectura objetivo (referencia)}

La siguiente arquitectura es una \textit{referencia orientativa}. Cada equipo puede proponer variaciones fundamentadas.

\begin{verbatim}
+---------------+         HTTP          +---------------+
|    React      | ===================> |    FastAPI    |
|  (Frontend)   | <================= |   (Backend)   |
+---------------+                       +---+---+---+--+
                                            |   |   |
                              Cache MISS    |   |   | Cache HIT
                                            v   |   v
                                      +-----+   | +------+
                                      | PostgreSQL| Redis |
                                      |         | |      |
                                      +---------+ +------+
                                                       ^
                                             Prometheus scrape
                                                       |
                                                 +-----+------+
                                                 |   Grafana  |
                                                 +------------+
\end{verbatim}

% ------------------------------------------------
\section{Descripción del Dominio de Negocio}
% ------------------------------------------------

\subsection{Aeronaves disponibles}

\begin{center}
\begin{tabular}{|l|c|l|}
\hline
\textbf{Modelo} & \textbf{Capacidad (pasajeros)} & \textbf{Operación} \\
\hline
Airbus A320       & 180 & Nacional \\
Boeing 737-800    & 189 & Nacional e internacional \\
Boeing 787 Dreamliner & 296 & Internacional \\
\hline
\end{tabular}
\end{center}

\subsection{Rutas}

\begin{center}
\begin{tabular}{|l|l|}
\hline
\textbf{Rutas nacionales} & \textbf{Rutas internacionales} \\
\hline
Santiago -- Antofagasta  & Santiago -- Buenos Aires \\
Santiago -- Concepción   & Santiago -- Lima \\
Santiago -- Puerto Montt & Santiago -- São Paulo \\
Santiago -- Punta Arenas & Santiago -- Miami \\
\hline
\end{tabular}
\end{center}

\subsection{Datos precargados (seed)}

El script \texttt{db/init.sql} pobla la base de datos con:

\begin{itemize}
  \item \textbf{3} registros de aeronaves.
  \item \textbf{8} rutas (4 nacionales + 4 internacionales).
  \item \textbf{200} vuelos programados con fechas de salida, precios y estado de ocupación.
  \item \textbf{500} reservas existentes distribuidas entre los vuelos.
\end{itemize}

% ------------------------------------------------
\section{Problema de Negocio}
% ------------------------------------------------

Durante un evento Cyber Day, miles de usuarios acceden simultáneamente al sistema para buscar vuelos y comprar tickets.

El endpoint crítico es:

\begin{lstlisting}[language=bash]
GET /flights
\end{lstlisting}

Este endpoint ejecuta la siguiente consulta (simplificada):

\begin{lstlisting}[language=sql, caption={Consulta del endpoint /flights (versión base)}]
SELECT
    f.id,
    f.departure_date,
    f.price,
    r.origin,
    r.destination,
    a.model AS aircraft,
    a.capacity,
    (a.capacity - COUNT(b.id)) AS available_seats
FROM flights f
JOIN routes r ON f.route_id = r.id
JOIN aircraft a ON f.aircraft_id = a.id
LEFT JOIN bookings b ON b.flight_id = f.id
GROUP BY f.id, r.id, a.id
ORDER BY f.departure_date;
\end{lstlisting}

\begin{infobox}[¿Por qué esto es un problema?]
Bajo 1.000 usuarios concurrentes, esta consulta se ejecuta centenares de veces por segundo. PostgreSQL no puede escalar linealmente con ese volumen de lecturas idénticas, generando aumento de latencia, timeouts y eventual indisponibilidad del servicio.
\end{infobox}

% ------------------------------------------------
\section{Simulación de Alta Concurrencia}
% ------------------------------------------------

El repositorio incluye un script base de Locust listo para ejecutar.

\subsection{Script base (referencia)}

\begin{lstlisting}[language=python, caption={locust/locustfile.py --- script base entregado}]
from locust import HttpUser, task, between

class SkyConnectUser(HttpUser):
    wait_time = between(1, 3)

    @task(5)
    def search_flights(self):
        self.client.get("/flights")

    @task(2)
    def get_routes(self):
        self.client.get("/routes")

    @task(1)
    def get_aircraft(self):
        self.client.get("/aircraft")
\end{lstlisting}

\begin{lstlisting}[language=bash, caption={Ejecutar prueba de carga}]
# Iniciar Locust (interfaz web en http://localhost:8089)
docker compose run --rm locust \
  locust -f locustfile.py \
  --host=http://backend:8000 \
  --users 1000 \
  --spawn-rate 50 \
  --run-time 5m \
  --headless \
  --csv=results/before
\end{lstlisting}

\subsection{Requisitos mínimos de las pruebas}

\begin{itemize}
  \item \textbf{Usuarios concurrentes:} mínimo 1.000.
  \item \textbf{Duración:} mínimo 5 minutos por ejecución.
  \item \textbf{Escenarios:} al menos dos ejecuciones comparativas (antes y después de optimizar).
  \item \textbf{Exportación:} resultados exportados en CSV y capturas de Grafana.
\end{itemize}

% ------------------------------------------------
\section{Requerimientos Obligatorios}
% ------------------------------------------------

\subsection{Implementación de Redis}

Cada equipo deberá integrar Redis como capa de caché entre FastAPI y PostgreSQL. La implementación debe considerar:

\begin{itemize}
  \item \textbf{Cache de vuelos} (\texttt{GET /flights}): resultado completo con TTL apropiado.
  \item \textbf{Cache de rutas} (\texttt{GET /routes}): datos relativamente estáticos.
  \item \textbf{Cache de aeronaves} (\texttt{GET /aircraft}): catálogo sin cambios frecuentes.
  \item \textbf{Estrategia de invalidación}: definir y documentar cuándo y cómo se invalida cada clave.
\end{itemize}

\begin{infobox}[Guía de implementación Redis]
\begin{enumerate}
  \item Agregar el servicio \texttt{redis} al \texttt{docker-compose.yml}.
  \item Instalar \texttt{redis-py} (o \texttt{aioredis}) en el backend.
  \item En cada endpoint, verificar primero la clave en Redis; si existe (\textit{cache hit}), retornar el valor; si no (\textit{cache miss}), consultar PostgreSQL, almacenar en Redis y retornar.
  \item Definir un TTL adecuado según la volatilidad del dato (p.\,ej., 60 segundos para vuelos, 300 para rutas).
  \item Exponer una métrica \texttt{cache\_hits\_total} y \texttt{cache\_misses\_total} para Prometheus.
\end{enumerate}
\end{infobox}

\subsection{Observabilidad}

Es obligatorio implementar Prometheus y Grafana. Los dashboards deben mostrar \textbf{como mínimo}:

\begin{center}
\begin{tabular}{|l|l|}
\hline
\textbf{Métrica}              & \textbf{Descripción} \\
\hline
Tiempo promedio de respuesta  & Latencia media del endpoint \texttt{/flights} \\
Latencia P95                  & Percentil 95 de tiempos de respuesta \\
Requests por segundo (RPS)    & Throughput total del sistema \\
Uso de CPU                    & Por contenedor (backend, db, redis) \\
Uso de memoria                & Por contenedor \\
Cache Hit Ratio               & \texttt{hits / (hits + misses)} \\
\hline
\end{tabular}
\end{center}

\begin{lstlisting}[language=bash, caption={Acceso a los servicios de observabilidad}]
# Prometheus:  http://localhost:9090
# Grafana:     http://localhost:3001  (admin / admin)
\end{lstlisting}

% ------------------------------------------------
\section{Mejoras Opcionales}
% ------------------------------------------------

Los equipos que implementen mejoras adicionales podrán alcanzar la categoría \textbf{Sobresaliente} en los criterios correspondientes. Se sugieren (pero no se limitan a):

\begin{itemize}
  \item Microservicio de búsqueda de vuelos desarrollado en \textbf{Go} (mayor throughput).
  \item \textbf{Índices} adicionales en PostgreSQL (p.\,ej., índice compuesto en \texttt{flights(route\_id, departure\_date)}).
  \item \textbf{Pool de conexiones} con PgBouncer para reducir la presión sobre PostgreSQL.
  \item \textbf{Balanceo de carga} con Nginx o Traefik para escalar horizontalmente el backend.
  \item \textbf{Cache de segundo nivel} con estrategia \textit{cache-aside} o \textit{read-through}.
  \item \textbf{Optimización de consultas SQL}: reescritura de la query con CTEs o vistas materializadas.
  \item \textbf{Rate limiting} para proteger el sistema ante picos anómalos.
  \item \textbf{Alertas en Grafana} configuradas con umbrales de SLA (p.\,ej., P95 > 500 ms).
\end{itemize}

\begin{warningbox}[Importante sobre las mejoras opcionales]
Cada mejora opcional debe estar respaldada por evidencia cuantitativa (métricas comparativas). No basta con implementar la mejora; es necesario demostrar su impacto con datos.
\end{warningbox}

% ------------------------------------------------
\section{Entregables}
% ------------------------------------------------

\subsection{Repositorio de código}

\begin{itemize}
  \item Repositorio Git con el código completo y funcional.
  \item Archivo \texttt{README.md} con instrucciones de despliegue reproducibles.
  \item \texttt{docker-compose.yml} con todos los servicios (backend, frontend, db, redis, prometheus, grafana, locust).
  \item Archivos de configuración de Prometheus (\texttt{prometheus.yml}) y dashboards de Grafana exportados en JSON.
  \item Resultados de las pruebas de carga en \texttt{/results} (CSV y capturas de pantalla).
\end{itemize}

\subsection{Documento técnico}

El informe debe incluir los siguientes apartados numerados:

\begin{enumerate}
  \item \textbf{Diagnóstico inicial:} descripción del sistema base y sus limitaciones observadas.
  \item \textbf{Identificación de cuellos de botella:} evidencia con métricas previas a las optimizaciones.
  \item \textbf{Arquitectura original:} diagrama y descripción del sistema recibido.
  \item \textbf{Arquitectura optimizada:} diagrama y justificación de cada componente añadido.
  \item \textbf{Mejoras implementadas:} detalle técnico de cada optimización (qué, cómo, por qué).
  \item \textbf{Métricas comparativas:} tablas y gráficas antes/después para cada mejora.
  \item \textbf{Justificación económica:} estimación de costos de infraestructura y beneficios obtenidos.
  \item \textbf{Evaluación esfuerzo vs.\ beneficio:} análisis crítico de qué mejoras aportaron más valor relativo al esfuerzo.
\end{enumerate}

\subsection{Diagramas requeridos}

\begin{itemize}
  \item Diagrama de casos de uso (flujo del usuario buscando y comprando un vuelo).
  \item Diagrama de componentes (arquitectura original).
  \item Diagrama de componentes (arquitectura optimizada).
  \item Diagrama de flujo de una solicitud \texttt{GET /flights} con Redis (cache hit y cache miss).
  \item Diagrama de secuencia del proceso de caché (opcional, para Sobresaliente).
\end{itemize}

\subsection{Presentación oral}

\begin{center}
\begin{tabular}{|l|c|}
\hline
\textbf{Instancia}     & \textbf{Tiempo} \\
\hline
Exposición              & 15 minutos \\
Preguntas del docente  & 5 minutos \\
\hline
\textbf{Total}          & \textbf{20 minutos} \\
\hline
\end{tabular}
\end{center}

La presentación debe incluir evidencia en vivo de:

\begin{itemize}
  \item Sistema corriendo con Docker Compose.
  \item Dashboards de Grafana con métricas activas.
  \item Prueba de carga en ejecución (o resultados exportados).
  \item Comparación antes/después con métricas concretas.
\end{itemize}

\newpage

% ------------------------------------------------
\section{Rúbrica de Evaluación}
% ------------------------------------------------

Cada criterio se evalúa en una escala de \textbf{1 a 4 puntos} según el nivel alcanzado.

\renewcommand{\arraystretch}{1.6}
\small

\begin{longtable}{|p{2.8cm}|p{2.9cm}|p{2.9cm}|p{2.9cm}|p{2.9cm}|}
\hline
\textbf{Criterio} &
\textbf{Insuficiente (1)} &
\textbf{Básico (2)} &
\textbf{Competente (3)} &
\textbf{Sobresaliente (4)} \\
\hline
\endfirsthead

\hline
\textbf{Criterio} &
\textbf{Insuficiente (1)} &
\textbf{Básico (2)} &
\textbf{Competente (3)} &
\textbf{Sobresaliente (4)} \\
\hline
\endhead

Diagnóstico del problema &
No identifica problemas relevantes ni presenta métricas iniciales. &
Identifica problemas de forma superficial sin evidencia cuantitativa. &
Identifica correctamente los cuellos de botella y los respalda con métricas. &
Análisis profundo, métricas objetivas, comparación con estándares de la industria. \\
\hline

Implementación de Redis &
No implementa Redis o no funciona en el entorno Docker. &
Implementación parcial; solo un endpoint usa caché o la invalidación es incorrecta. &
Redis implementado y funcional con TTL y estrategia de invalidación documentados. &
Redis optimizado: múltiples estrategias de caché, métricas de hit/miss, análisis de TTL óptimo. \\
\hline

Observabilidad (Prometheus + Grafana) &
No existen métricas ni dashboards operativos. &
Dashboards presentes pero incompletos; faltan métricas críticas. &
Dashboards funcionales con todas las métricas mínimas requeridas. &
Observabilidad completa con alertas configuradas, análisis de tendencias y correlación de métricas. \\
\hline

Pruebas de carga (Locust) &
No realiza pruebas o no documenta los parámetros utilizados. &
Pruebas ejecutadas con menos usuarios o tiempo del mínimo requerido. &
Pruebas correctamente ejecutadas ($\geq$1.000 usuarios, $\geq$5 min) y documentadas. &
Pruebas rigurosas con múltiples escenarios, análisis estadístico y comparación detallada. \\
\hline

Métricas comparativas &
No presenta comparación entre sistema base y sistema optimizado. &
Comparación parcial o sin datos suficientes para sacar conclusiones. &
Compara adecuadamente las métricas clave antes y después de cada optimización. &
Comparación exhaustiva con gráficas, tablas y análisis estadístico por mejora implementada. \\
\hline

Justificación económica &
No existe análisis económico. &
Análisis superficial sin cuantificación de costos ni beneficios. &
Estimación razonada de costos de infraestructura y ahorro obtenido. &
Análisis sólido con proyección financiera, cálculo de ROI y comparación de alternativas. \\
\hline

Esfuerzo vs.\ beneficio &
No analiza la relación entre el esfuerzo invertido y los resultados. &
Menciona las mejoras implementadas sin análisis crítico. &
Analiza adecuadamente qué mejoras fueron más rentables en relación al esfuerzo. &
Evaluación estratégica con priorización fundamentada y recomendaciones para trabajo futuro. \\
\hline

Diagramas y documentación &
Documentación incompleta, inconsistente o sin diagramas. &
Diagramas básicos presentes pero sin suficiente detalle técnico. &
Documentación clara, completa y con todos los diagramas requeridos. &
Documentación de calidad profesional con diagramas detallados y README reproducible. \\
\hline

Presentación oral &
No comunica adecuadamente los resultados; no responde preguntas. &
Presentación parcialmente estructurada; dificultad para responder preguntas técnicas. &
Presentación clara, técnicamente correcta y con evidencia en vivo. &
Presentación profesional, dominio total de la solución, respuestas precisas y fundamentadas. \\
\hline

\end{longtable}

\normalsize

% ------------------------------------------------
\section{Ponderación}
% ------------------------------------------------

\begin{center}
\begin{tabular}{|l|c|c|}
\hline
\textbf{Criterio} & \textbf{Peso} & \textbf{Puntos máx.} \\
\hline
Diagnóstico del problema        & 10\% & 4 \\
Implementación de Redis         & 20\% & 4 \\
Observabilidad y métricas       & 15\% & 4 \\
Pruebas de carga (Locust)       & 10\% & 4 \\
Métricas comparativas           & 10\% & 4 \\
Justificación económica         & 10\% & 4 \\
Evaluación esfuerzo--beneficio  & 5\%  & 4 \\
Diagramas y documentación       & 10\% & 4 \\
Presentación oral               & 10\% & 4 \\
\hline
\textbf{Total}                  & \textbf{100\%} & --- \\
\hline
\end{tabular}
\end{center}

\vspace{0.4cm}

\textbf{Fórmula de la nota:}

\[
  \text{Nota} = \sum_{i=1}^{9} \left( \frac{\text{Puntos obtenidos}_i}{4} \times \text{Peso}_i \right) \times 6 + 1
\]

Donde 1 corresponde a la nota mínima y 7 a la nota máxima (escala chilena 1--7).

% ------------------------------------------------
\section{Criterio de Aprobación}
% ------------------------------------------------

El hito se considerará aprobado cuando se cumplan \textbf{todos} los siguientes requisitos:

\begin{itemize}
  \item El sistema corre completamente con \texttt{docker compose up} sin errores manuales.
  \item Existe una reducción medible y documentada de la latencia en el endpoint \texttt{GET /flights}.
  \item Redis está efectivamente integrado y el Cache Hit Ratio supera el 70\% durante las pruebas de carga.
  \item Los dashboards de Grafana muestran métricas en tiempo real durante la presentación.
  \item Las mejoras están respaldadas por evidencia cuantitativa (CSV de Locust + capturas de Grafana).
  \item Se presenta una justificación económica con valores estimados fundamentados.
  \item La documentación y el repositorio son suficientemente claros para que otro equipo reproduzca el entorno.
\end{itemize}

% ------------------------------------------------
\section{Consideraciones Éticas y de Integridad Académica}
% ------------------------------------------------

\begin{itemize}
  \item Todo el código debe ser desarrollado por el equipo. El uso de herramientas de asistencia de IA (p.\,ej., GitHub Copilot, ChatGPT) está permitido, pero \textbf{debe declararse} en el informe especificando qué partes fueron asistidas y cómo fueron revisadas.
  \item Los resultados de las pruebas de carga no pueden ser fabricados o editados manualmente.
  \item Los integrantes del equipo deben ser capaces de explicar cualquier parte del código durante la presentación.
  \item La detección de plagio o datos fabricados implica calificación inmediata con nota 1 en el hito.
\end{itemize}

% ------------------------------------------------
\section{Preguntas Frecuentes}
% ------------------------------------------------

\begin{description}
  \item[¿Puedo cambiar el stack del frontend?] No es obligatorio ni recomendado; el foco está en el backend y la capa de caché.

  \item[¿Puedo agregar más endpoints?] Sí, pero los criterios de evaluación se centran en optimizar \texttt{GET /flights}. Otros endpoints son opcionales.

  \item[¿Qué pasa si Redis no mejora la latencia?] Es un resultado válido si está bien documentado y analizado. Lo importante es la rigurosidad del proceso, no solo el resultado positivo.

  \item[¿El TTL de Redis es fijo?] No. El equipo debe justificar el TTL escogido según la volatilidad del dato y el impacto en la consistencia.

  \item[¿Debo implementar autenticación?] No. El sistema base no incluye autenticación y no es requerida para este hito.
\end{description}

\end{document}
